Modern attention-based machine learning (ML) models
are essential for language, vision, and other emerging tasks in artificial intelligence
\cite{gpt5,image-gen-paper,audio-gen-paper}. To efficiently deploy these models on GPUs,
system engineers need tools to bridge the performance gap
between a \emph{math-like model specification}, such as PyTorch or JAX code, and efficient
low-level code that maximally utilizes the underlying hardware.
Today, \NA{existing tools fall short of this goal on increasingly varied hardware:}
they are either too hard to use, requiring expertise in rapidly-evolving low-level
hardware details, or produce programs that are not sufficiently performant.

The most common approach today is writing manually optimized kernels in hardware vendor languages
like CUDA~\cite{cuda} and CUTLASS~\cite{cutlass}. They achieve high performance and have been
used in PyTorch~\cite{pytorch} and TensorFlow~\cite{tensorflow}.
However, programming in these low-level languages requires significant hardware and
algorithmic expertise and does not scale to rapidly evolving GPU architectures and ML models.
On the opposite side, schedule-based tensor compilers such as Halide~\cite{halide} and
TVM~\cite{tvm} optimize a high-level specification of the ML operator {in a math-like
tensor language.}
This design offloads the obligation of optimization from the programmer to the compiler,
which needs to search in a complex design space of \emph{schedules}
(sequence of transformations) to find an optimal kernel on the target GPU.
\NA{But this space mixes high-level decisions (e.g., fusion) with low-level ones
(e.g., thread mappings)}, and searching it is intractable; thus
today's high-level tensor compilers fail to match handwritten code in performance.

As a tentative trade-off between programmability and efficiency,
tile-based compilers and languages,
including Triton~\cite{triton}, Pallas~\cite{pallas}, and TileLang~\cite{tilelang},
have emerged as an abstraction layer over hardware vendor languages.
Here, a programmer manually expresses the ML operator as computation over tiles (fixed
size subtensors), and the tile compiler applies low-level SIMD
optimizations to map tile computation to the hardware.
With tile languages, a programmer still has to manually control many essential
high-level optimizations. This issue is especially evident for attention:
each version of FlashAttention
\cite{flash-attention,flash-attention-2,flash-decoding,flash-attention-3}
is an increasingly complex, manually written tile program
that must be significantly modified for every new GPU architecture.

\NA{We advocate for a \emph{math-to-kernel} compiler}:
the program\-mer writes operators in a math-like linear algebraic language
(e.g., attention as a sequence of matmuls and softmax),
and the compiler automatically finds high-performance kernels (e.g., FlashAttention).
\NA{Such a compiler is a viable solution that can yield ML
kernel efficiency competitive with manually-written code on modern GPU architectures.}
The central difficulty to materialize this goal is that we must bridge two types of optimizations
that live on different levels of abstraction:
(1)~operation-level optimizations that reason about compute expressions and data
dependencies of linear algebra operations,
and (2)~tile optimizations that map regular computation patterns to the underlying hardware.
\NA{Our key insight is that the two levels can be separated:
  tile compilers already perform (2) effectively and can serve as the compiler's
  low-level IR and backend, so schedules only need to capture (1),
  which drastically shrinks the space to search.
Automating (1) on top of tile compilers, however, raises two challenges.}
First, compiling many tensor operations into one or a few efficient kernels
requires new fusion techniques to handle complex data dependency.
In particular, \emph{advanced reduction fusion}, introduced by Neptune~\cite{neptune},
can fuse reduction loop nests, such as softmax and matmul in attention,
and it is a key optimization to allow compilers to automatically discover
FlashAttention from math-like attention.
\NA{Neptune already delegates tile optimizations to Triton,
but it relies on manually written schedules;}
it is still unclear how to apply such optimizations in a compiler without user intervention.
\NA{Second, no single intermediate representation (IR) suits all of these optimizations.
  Operation-level optimizations reason about both data dependencies (e.g., fusion)
  and compute expressions (e.g., expression rewriting):
  fusion needs loops and data dependencies,
  while expression rewriting needs tiles as values it can freely manipulate.
Tile compilers, in turn, need explicit tile memory accesses.}
% \sasa{We need to develop the argument more here. Why is tile promising and why is tile
% not a good language for programmers but an IR?}

\vspace{0.3em}
\noindent{\textbf{Our work: \NAME{}}. } \NAME{} is a novel tensor compiler for
attention-based ML models.
\NAME{} brings together operation-level optimizations
(advanced reduction fusion and expression rewriting)
and tile optimizations in one pipeline, to generate high-performance
GPU kernels on modern and emerging GPUs.
\NAME{} adopts the concept of schedules to regulate the legality of the sequence of transformations
and facilitate search-based exploration~\cite{tvm},
and automates the generation of schedules that incorporate advanced reduction fusion
via a novel auto-scheduler designed specifically to target tile compilers.
The auto-scheduler addresses three challenges that have kept such optimizations manual.
\NA{(1) \emph{Non-local fusion:} advanced fusion fuses two reductions together with all
  element-wise operations between them,
  so a pass that transforms one statement at a time cannot decide it locally.
  The auto-scheduler looks ahead along data dependencies to find these opportunities.
  (2) \emph{Tile-readiness:} operation-level transformations can leave the program in a
  form tile compilers handle poorly, e.g., advanced fusion can introduce scan dependencies.
  The auto-scheduler regularizes the program to keep it ``tile-ready'',
  i.e., easy to optimize by tile compilers.
  (3) \emph{Input-dependent trade-offs:} the best schedule to apply and tile compiler to use
  depend on the operator, input shape, and hardware.
The auto-scheduler generates promising schedules and lets the auto-tuner select among them.}

\NA{\NAME{} then successively lowers the scheduled program through three IRs,
  each at the abstraction level of one category of optimizations.
  Fusion and other scheduling happen in a scalar IR.
  Expression rewriting happens in a novel virtual-register tile IR (VR-tile IR),
  which expresses tiles as values rather than memory accesses.
  Tile optimizations happen in a memory-access tile IR (MA-tile IR),
  which \NAME{} translates to multiple tile compilers (Triton, Tawa~\cite{tawa}, and TileLang)
and selects the best-performing one.}

\NAME{} is the first end-to-end compiler to perform fully automated math-to-kernel
optimization on attention: \NA{starting from a simple linear algebraic representation of attention,
it automatically discovers variants of FlashAttention-3 and delivers FlashAttention-3 performance.}
%\sasa{be careful re Flashlight}
%\sasa{This is the first general end-to-end compiler system that can do this task. (we
% may need to massage it!)}

We evaluate \NAME{} on five \matei{needs updating}
% \sasa{[Is it five?]}
transformer-based models, with {150} total evaluation configurations, on two
modern GPU architectures: NVIDIA GH200 (Hopper) and NVIDIA RTX 5090 (Blackwell). We used
seven compiler backends and five hand-optimized libraries as baselines. \NAME{}-generated
attention kernels have up to 23\% higher throughput over the
state-of-the-art compiler systems on NVIDIA GH200 and up to 42\% on NVIDIA RTX 5090.
Moreover, our results are comparable to and for many configurations (with longer sequences)
faster than manually written cuDNN kernels.

\vspace{0.3em}
\noindent\textbf{Contributions:} The paper makes several contributions:
\begin{itemize}
  \item \NAME{}, a tensor compiler for attention-based ML models that fully automates the
    generation of high-per\-for\-mance GPU kernels from a program description (\S\ref{sec:design}).
  \item An auto-scheduler composing advanced reduction fusion and other high-level
    optimizations in \NAME{}'s scheduling language, while keeping programs tile-ready
    (\S\ref{sec:auto-sched}).
  \item A three-IR lowering pipeline: a scalar IR, a novel VR-tile IR for expression
    rewriting, and an MA-tile IR that enables multiple tile compilers as backends
    (\S\ref{sec:opt-pipeline}).
  \item An evaluation on multiple state-of-the-art attention models and GPU architectures,
    showing higher throughput than existing compilers and performance on par with
    hand-optimized kernels (\S\ref{sec:eval}).
\end{itemize}
